% !TeX encoding = UTF-8 !TeX spellcheck = en_US
\documentclass[11pt, a4paper]{article}

% The title of the current document to be produced.
\newcommand{\doctitle}{Assignment \#7}
\newcommand{\assignments}{\bluetext{\textbf{Assignment:}}}
%
\setlength{\unitlength}{1in}
\renewcommand{\arraystretch}{2}
\setlength{\columnsep}{2cm}

\newcommand\dunderline[3][-1pt]{{%
  \sbox0{#3}%
  \ooalign{\copy0\cr\rule[\dimexpr#1-#2\relax]{\wd0}{#2}}}}

\usepackage[sfdefault]{atkinson}
\usepackage[T1]{fontenc}

\def\code#1{\bluetext{\texttt{#1}}}


%------------------------------------------------------------
% Import commands for both teacher and course information.  | NOTE: Change your
% teacher and course info in these files. |
%------>------>------>------>------>------>------>------>-->|
%
%------------------------------------------------------------
%-- Import packages and custom command definitons.          |
%------>------>------>------>------>------>------>------>-->|
%% ==================================
%% ===== Teacher info             ===
%% ==================================
%%
\newcommand{\instructor}{Katelyn Breivik}
\newcommand{\office}{Wean 8319}
\newcommand{\hours}{3pm on Mondays}
\newcommand{\phone}{}
\newcommand{\college}{}
\newcommand{\email}{kbreivik@andrew.cmu.edu}
\newcommand{\faculty}{}
\newcommand{\department}{Department of Physics}
                              %|
%% ==================================
%% ===== Course-specific commands ===
%% ==================================

%- Instructions: change course info here. 
\newcommand{\semester}{Fall 2023}

\newcommand{\ponderation}{2-4-3 (Theory-Lab-Homework)}
\newcommand{\coursetitle}{Astrophysics of Stars and the Galaxy}
\newcommand{\coursenumber}{[33-467]}
\newcommand{\prerequisite}{33-224, 33-234}
 
\input{../includes/packages-imports}                          %|  
\input{../includes/custom-commands}   
%
%---> Genereate & inject metadata                           |
\input{../includes/hyperef.doc.info}                          %|
%------------------------------------------------------------

\topmargin -70pt
\begin{document} 

\normalfont

%-------------------------------------------------------------
%-- Make the header of the document                          |
%------>------>------>------>------>------>------>------>--> |
%---------------------------------------------------------------------
%-- The following produces the document header including the title.  |
%---------------------------------------------------------------------

\noindent % <-- need to have this first.

\hfill	
\begin{minipage}{0.60\textwidth}%
	\raggedleft%
	{\Large \textsc{\doctitle}\par}
	\doublerule 
\end{minipage}%
\vspace{0.9cm}

  
%-------------------------------------------------------------
%-- Problem # 1                                              |
%------>------>------>------>------>------>------>------>--> |
\noindent \dunderline{1.0pt}{\textbf{Problem \#1a: Accessing public exopolanet
data (5 pts)}}\\
The largest catalog of exoplanets can be found in the
\href{https://exoplanetarchive.ipac.caltech.edu/cgi-bin/TblView/nph-tblView?app=ExoTbls&config=PSCompPars}{Composite
Data Table} of the NASA Exoplanet Archive. Download the data table by clicking
on the "Download" button, and save to a datafile that you can open locally on
your computer/laptop. To verify that you have accessed the data, plot a
histogram of the log-base-10 of the orbital period of each exoplanet. Be sure to
label your axes!

\vspace{4pt}
\noindent\emph{Hint:} there are several header columns in the data file that
will be dowloaded. These headers show you the names of the columns (including
the name of the orbital period column).

\vspace{4pt}
\noindent\emph{Hint 2:} Pandas makes loading data easy. You can specify the
number of header lines to skip using `skiprows=number\_to\_skip' in the function
call that loads the data. If you skip the right number of rows, you can use the
row just before the data stars to specify the columns automatically!

\vspace{8pt}
\noindent \dunderline{1.0pt}{\textbf{Problem \#1b: Multiplanet systems (10
pts)}}\\
Now that you have data loaded, you should use the column which specifies the
number of planets per system to plot histograms for subsets of the data
containing 1, 2, 3, 4, and 5 planets per system on the same figure. Be sure to
include a legend so that you can see how the orbital period distributions
change.

\vspace{8pt}
\noindent \dunderline{1.0pt}{\textbf{Problem \#1c: Biases in observing
multiplanet systems (9 pts)}}\\
Based on your figure and the material in the class lectures, provide two
examples of how observational biases could affect the size and
orbital period distribution of planetary systems. Provide one argument, based on
exoplanet dynamics, for why there could truly be fewer multiplanet systems.

\vspace{20pt}

\noindent \dunderline{1.0pt}{\textbf{Problem \#2: Figures and captions (again!)
(20 points)}}\\
\noindent Your project presentations are coming up soon! As part of your
presentation, you should plan to include at least one original figure that you
(and your partner if you have one) make. In the presentation, you'll be able to
walk the audience through the figure. This isn't the only way that figures are
presented though!\\

\noindent For this problem, you should create a figure (you and your partner can
submit the same figure) that you could use in your presentation. The figure
should illustrate some conclusion or result from your project and should have
adequate axis labels and a legend if necessary. Once you have a figure, you
should also write a figure caption that describes the necessary information for
a reader to understand and interpret the figure themselves. It's also a good
idea to add in a punchline that you want your reader to take away from the
figure. \emph{You must write your own figure caption even if you and your
partner share the figure.}\\

\end{document} 
